\documentclass{article}
\usepackage[UTF8]{ctex}
\usepackage{amsmath,amssymb,amsfonts}
\usepackage{enumitem}
\usepackage{geometry}
\geometry{a4paper, margin=1in}
% === 错题本排版优化（不参与题目标识） ===
\setlength{\parindent}{0pt}
\setlength{\parskip}{0.35em plus 0.10em minus 0.08em}
\setlength{\emergencystretch}{3em}
\setlist{itemsep=0.25em, topsep=0.40em, parsep=0pt}
\allowdisplaybreaks[2]
% === 错题本排版优化结束 ===
\begin{document}
\section*{东北师范大学2025年数学分析真题}
\section*{79．东北师范大学2025年数学分析试卷}

\subsection{题目1}
1． \(\lim _{n \rightarrow \infty}\left(1+\frac{1}{2}+\cdots+\frac{1}{n}\right)^{\frac{1}{n}}\) ．

\subsection{题目2}
2． \(\lim _{x \rightarrow 0} \frac{\int_{0}^{\sin ^{2} x} t e^{t} \mathrm{~d} t}{\sqrt{1+x^{4}}-1}\) ．

\subsection{题目3}
3． \(\lim _{(x, y) \rightarrow(\infty, \infty)} \frac{x+y}{x^{2}-x y+y^{2}}\) ．

\subsection{题目4}
4． \(\iint_{D}\left[x^{2}+y^{2}\right] \mathrm{d} x \mathrm{~d} y\) ，其中 \(D=\{(x, y) \mid-1 \leq x \leq 1,-1 \leq y \leq 1\}\) ，函数［•］是取整函数．

\subsection{题目5}
5． \(\int_{-\pi}^{\pi}\left(2+\sum_{n=1}^{100} \frac{\cos n x}{n}\right)^{2} \mathrm{~d} x\) ．

二．（15 分）计算曲线积分

\[
\int_{C}\left(4 x^{3} y^{3}-3 y^{2}+6 y+1\right) \mathrm{d} x+\left(3 x^{4} y^{2}-6 x y+10 x-4\right) \mathrm{d} y
\]

积分路径 \(C\) 是圆弧 \(x^{2}+y^{2}=1(x \geq 0, y \geq 0)\) ，起点是 \((1,0)\) ，终点是 \((0,1)\) 。

三．（15 分）设 \(u=e^{x+y}\) ，而 \(y\) 是由方程 \(y-\frac{1}{2} \sin y=x\) 所确定的函数，求 \(\frac{\mathrm{d} u}{\mathrm{~d} x}, \frac{\mathrm{~d}^{2} u}{\mathrm{~d} x^{2}}\) ．

四．（15 分）设 \(f(x)\) 在 \([a, b]\) 上连续，在 \((a, b)\) 上可导，\(f(a)=0\) ，当 \(x \in(a, b]\) 时，\(f(x)>0\) ．证明： \(\frac{f^{\prime}(x)}{f(x)}\) 在 \((a, b)\) 内无界．

五．（15 分）设 \(a>b>0\) ，求级数 \(\sum_{n=1}^{\infty} \frac{a^{n}+(-b)^{n}}{2 n}(x-1)^{n}\) 的和函数．

六．（15分）设函数 \(f(x)\) 在区间 \((-\infty,+\infty)\) 内连续可导，令

\[
f_{n}(x)=\frac{n}{2}\left[f\left(x+\frac{1}{n}\right)-f\left(x-\frac{1}{n}\right)\right], n=1,2,3, \cdots .
\]

（1）证明：对任意的闭区间 \([\alpha, \beta]\) ，函数列 \(\left\{f_{n}(x)\right\}\) 在 \([\alpha, \beta]\) 上一致收敛．\\
（2）求 \(\lim _{n \rightarrow \infty} \int_{0}^{2024} f_{n}(x) \mathrm{d} x\) ．

七．（15 分）设数项级数 \(\sum_{n=1}^{\infty} a_{n}\) 收敛，\(\sum_{n=1}^{\infty}\left(b_{n+1}-b_{n}\right)\) 绝对收敛，证明：\(\sum_{n=1}^{\infty} a_{n} b_{n}\) 收敛。

八．（15 分）设 \(a_{0} \in(-\infty,+\infty), a_{n}=\cos a_{n-1}(n=1,2,3, \cdots)\) ，证明：数列 \(\left\{a_{n}\right\}\) 收敛．

九．（15 分）设函数 \(f(x)\) 在 \([0,1]\) 上单调增加，\(f(0)>0, f(1)<1\) ．证明：对任意的 \(\alpha>0\) ，存在\(\xi \in(0,1)\) ，使得 \(f(\xi)=\xi^{\alpha}\) 。
\end{document}
